\chapter{ServiceManager}
\label{SrvcMgr}

This chapter provides a guidelines about the usage of the API configured in the Service Manager.

\section{API SrvcMgr\_ActivateServices}
This API is used to set a specific state for a list of Service ID. The result of these operations are retrieved to the caller of the API in the commandStatusList.
The number of Service ID (length of serviceIdList) must be specified in the listLength by the caller of the API and must match with the length of the 
activationCommandList and commandStatusList. 
The same index for each of these three arrays correspond a Service ID on which we want to set a specific state and the result in retrieved in the same index
of the commandStatusList.

Each member of the array activationCommandList can assume just two value (TRUE or FALSE):
\begin{itemize}
\item TRUE: to set a Server Service to SD\_SERVER\_SERVICE\_AVAILABLE or a Client Service SD\_CLIENT\_ SERVICE\_REQUESTED (unformally to enable a Service);
\item FALSE: to set a Server Service to SD\_SERVER\_SERVICE\_DOWN or a Client Service SD\_CLIENT\_ SERVICE\_RELEASED (unformally to disable a Service); 
\end{itemize}
Please note that this API selects the Sd API based on the Service Group configuration. If the service belongs to a Service Group, SrvcMgr uses Sd\_ServiceGroupStart()/Sd\_ServiceGroupStop() and does not call Sd\_ServerServiceSetState() or Sd\_ClientServiceSetState(). If the service does not belong to a Service Group, SrvcMgr uses Sd\_ServerServiceSetState() for server services and Sd\_ClientServiceSetState() for client services.
For non-group services, the helper distinguishes client vs server instances, maps the activation flag to Sd\_ClientServiceSetState() or Sd\_ServerServiceSetState(), and updates commandStatusList with E\_OK\_REMOTE\_ENABLED or E\_OK\_LOCAL\_ENABLED on E\_OK. When a client service is requested and owns consumed event groups, Sd\_ConsumedEventGroupSetState() is called for each configured group; any E\_NOT\_OK moves the related status to E\_GENERAL\_ERROR. The same helper is reused by Sum\_Srvcmgr\_SetClientSrvcRequestedRemoteSrvcAvailable() so that service groups and non-group services share consistent activation behaviour.
Refer to the Sd module specification to understand more deeply the meaning of SD\_SERVER\_SERVICE\_AVAILABLE, SD\_CLIENT\_SERVICE\_REQUESTED, SD\_SERVER\_ SERVICE\_DOWN and SD\_CLIENT\_SERVICE\_RELEASED.


Each member of the array commandStatusList can assume the following value:
\begin{itemize}
\item E\_OK\_LOCAL\_ENABLED (0x00): Service is calibrated On and is a local service (Sd Server Service); 
\item E\_OK\_REMOTE\_ENABLED (0x01): Service is calibrated On and is a remote service (Sd Client Service);
\item E\_DISABLED (0x02): Service is calibrated Off and is local or remote Service (Sd Server or Sd Client Service);
\item E\_GENERAL\_ERROR (0x03): Service not configured or the SD APIs used returns E\_NOT\_OK for that particular Server or Client Service; 
\end{itemize}

Moreover, this API is triggered by the C/S interface SrvcMgr\_SPM\_ActivateServices. 
